\subsection{The next cyclotomic base change and its actual class quotients}

\begin{statement}{E.94}{Proved}{The \(\mu_{27}\) comparison, a noncyclic S-class quotient, and its exact norm and extension maps}
Keep the literal roots and primes of E.86 and E.91. Put
\[
 \eta^3=\xi,\qquad
 E=B_3(\eta),\qquad F=B_1,\qquad
 \beta_3^9=\beta_1,\qquad \xi^3=\zeta .
 \tag{E94.1}
\]
Then \([E:F]=27\), \([E:B_3]=3\), and
\(\operatorname{Gal}(E/F)=C_9\oplus C_3\).
At each of the three primes \(Q_i\) over three,
\(E/B_3\) is totally ramified. At each of the three
primes \(T_i\) over seventeen it is unramified of
residue degree three. All six primes have a unique
prolongation, denoted \(Q_i^E,T_i^E\). The ideal
norm and extension coefficients are therefore
\[
 \begin{aligned}
 N_{E/B_3}Q_i^E&=Q_i,&j_{E/B_3}Q_i&=3Q_i^E,\\
 N_{E/B_3}T_i^E&=3T_i,&j_{E/B_3}T_i&=T_i^E .
 \end{aligned}
 \tag{E94.2}
\]
Both ordinary and S-class norms onto \(B_3\) are
surjective, where \(S=\{3,17\}\).

There are literal \(K\)-automorphisms
\[
 \begin{aligned}
 G(\beta_3)&=\eta\beta_3,&G(\eta)&=\eta,& |G|&=27,\\
 H(\beta_3)&=\beta_3,&H(\eta)&=\eta^4,& |H|&=9,\\
 HGH^{-1}&=G^4,&
 \operatorname{Gal}(E/B_3)&=\langle H^3\rangle .
 \end{aligned}
 \tag{E94.3}
\]
They generate \(\operatorname{Gal}(E/K)\), of order
\(243\). The restrictions of \(G^3,H\) to \(B_3\)
are exactly \(g_2,h_2\) of E.91. The element \(G\)
does not preserve \(B_3\); indeed
\(GH^3G^{-1}=G^{-9}H^3\).

Let \(\mathcal G_E/F\) be the maximal abelian
three-extension containing \(E\) that is unramified
over \(E\). Let \(\mathcal G_E^S\) be its maximal
subextension in which all six S-primes of \(E\)
split. These fields and their actual restriction
characters are computed by the ray presentations
below. Their groups are
\[
 \begin{array}{c|c|c}
 &\operatorname{Gal}(\mathcal G/F)
      &\operatorname{Gal}(\mathcal G/E)\\ \hline
 \mathcal G_E&C_9^4\oplus C_3^2&C_9^3\oplus C_3\\
 \mathcal G_E^S&C_9^2\oplus C_3^2&C_9\oplus C_3 .
 \end{array}
 \tag{E94.4}
\]
Write the right-hand groups as
\(\mathcal A=C_9^3\oplus C_3\) and
\(\mathcal A_S=C_9\oplus C_3\). Restriction of
the Hilbert class fields gives actual surjections
\[
 \operatorname{Cl}(E)\{3\}\longrightarrow\mathcal A,
 \qquad
 \operatorname{Cl}_S(E)\{3\}\longrightarrow\mathcal A_S .
 \tag{E94.5}
\]
These are evaluated class quotients. The statement
does not assign either full class group to its quotient.

The bases in (E94.4) are specified by the integer
Smith maps in the certificate. In these bases, removal
of the S-primes is the exact matrix
\[
 Q:\mathcal A\longrightarrow\mathcal A_S,\qquad
 Q=\begin{pmatrix}
 0&8&3&6\\
 2&2&1&0
 \end{pmatrix}.
 \tag{E94.6}
\]
Its kernel, of order \(81\), is generated by the
following six actual prime classes:
\[
 C_E=
 \begin{array}{c|rrrrrr}
 &Q_1^E&Q_2^E&Q_3^E&T_1^E&T_2^E&T_3^E\\ \hline
 C_9&1&7&1&3&3&3\\
 C_9&3&6&0&0&0&0\\
 C_9&1&1&7&6&6&6\\
 C_3&0&2&1&0&0&0
 \end{array}.
 \tag{E94.7}
\]
The first three columns have order nine and the
last three have order three in this genus quotient.
Their images are computed from local norm
uniformizers, with the residue-degree coefficients
in (E94.2) retained.

Let \(A_9=C_9^2\oplus C_3\) have the simple
\((y,z,t)\) coordinates of E.91. The original
ninth-genus field embeds in \(\mathcal G_E\).
The resulting norm and extension comparisons are
\[
 \begin{gathered}
 N:\mathcal A\twoheadrightarrow A_9,\qquad
 N=\begin{pmatrix}
 3&1&0&3\\
 7&0&2&0\\
 0&0&2&0
 \end{pmatrix},\\[2mm]
 J:A_9\longrightarrow\mathcal A,\qquad
 J=\begin{pmatrix}
 0&3&6\\
 3&0&0\\
 0&0&6\\
 0&0&0
 \end{pmatrix},\qquad
 JN=3,\quad NJ=3,\quad \ker N=C_3^2 .
 \end{gathered}
 \tag{E94.8}
\]
The map \(N\) is induced by actual ideal norm.
The map \(J\) is induced by actual ideal extension:
the E.91 residual
\[
 W=(g_2-1)\operatorname{Cl}(B_3)\{3\}
 \quad\hbox{maps to zero in }\mathcal A .
 \tag{E94.9}
\]
Thus the extension comparison factors through
the already evaluated quotient \(A_9\); the proof
identifies that factor rather than choosing an
unrelated homomorphism between groups of the same order.

The normalizer action on the relative S-genus
quotient, and its actual norm to the original
S-class group, are
\[
 \begin{gathered}
 G|_{\mathcal A_S}=I,\qquad
 H|_{\mathcal A_S}=\begin{pmatrix}1&3\\2&1\end{pmatrix},
 \qquad
 \nu:\mathcal A_S\twoheadrightarrow
      \operatorname{Cl}_S(B_3)\{3\}=C_3,\\
 \nu(x,y)=-x\pmod3,\qquad \ker\nu=C_3^2 .
 \end{gathered}
 \tag{E94.10}
\]
The norm is identified using the coherent prime-\(73\)
basis of E.86. In particular
\(\operatorname{Cl}_S(E)/(H^3-1)\operatorname{Cl}_S(E)\)
has the noncyclic quotient \(C_9\oplus C_3\).
Consequently \(\operatorname{Cl}_S(E)\{3\}\) is not
a cyclic module over
\(\mathbb Z_3[\operatorname{Gal}(E/B_3)]\).
This follows through (E94.5), on which \(H^3\)
acts identically.

The computation uses the same certified degree-eighteen
field \(F\) as E.76. It introduces no BNF calculation
in \(B_3\), \(E\), or the original degree-eighteen
field \(L_2\). Every local kernel, ray relation,
Kummer character, Smith transformation and inverse
is retained. It also records the unsuccessful initial
local-unit spanning family and the corrected literal
norm family. The full class groups of \(E\), the
value of the E.91 transvection, and the eventual
exponent \(h\) of E.69 remain outside the evaluations
claimed here; equation (E94.9) states the exact map
by which that transvection disappears in this
particular comparison.

The manuscript anchor is section51.16, folios733v--736v,
scans172--178, via the actual positional norm systems
of E.46, E.69, E.72 and E.88.
\end{statement}
\begin{proof}
\emph{The fields, local degrees, and actual normalizer.}
The residue field at each prime above seventeen in
\(B_3\) is \(\mathbb F_{289}\). It contains \(\mu_9\)
and does not contain \(\mu_{27}\), since
\(v_3(289-1)=2\). Adjoining \(\eta\) is therefore
unramified of degree three there. This proves the
global degree \([E:B_3]=3\), and proves the assertion
at seventeen in (E94.2).

At a prime of \(F\) above three the completion is
\(F_v=\mathbb Q_3(\mu_9)\). Indeed \(17\) has a
cube root \(a_0\in\mathbb Q_3\), because
\(-17\in1+9\mathbb Z_3\), and the three embeddings
send \(a\) to \(a_0\zeta^j\). The local radical
extension \(B_3/F\) is cyclic of degree nine,
generated by a ninth root of \(\beta_1\).
Its local norm character will be denoted
\(\chi_v:F_v^\times\to\mathbb Z/9\).
The extension \(F_v(\eta)/F_v\) has its cyclotomic
character \(\psi_v:F_v^\times\to\mathbb Z/3\),
where the latter coordinate means
\(\eta\mapsto\zeta\eta\).

Their normalizations can be checked directly.
The actual arithmetic local Artin map on
\(\mathbb Q_3\)-units acts on \(\mu_{27}\) by
the inverse unit. This follows also by applying
global reciprocity to a rational element and
the prime Frobenii in the cyclotomic field
\cite[V,Example3.2,Proposition5.2,Theorem5.3 and
Summary5.10; III,Theorem3.4]{E_MilneCFT}.
Local norm compatibility then gives
\[
 \begin{gathered}
 N_{F_v/\mathbb Q_3}(\xi)=1,\qquad
 N_{F_v/\mathbb Q_3}(\beta_1)=17^2\equiv19\pmod{27},\\
 N_{F_v/\mathbb Q_3}(1+\beta_1)
       =1-17+17^2=273,\\
 \psi_v(\xi)=0,\qquad
 \psi_v(\beta_1)=1,\qquad
 \psi_v(1+\beta_1)=2.
 \end{gathered}
 \tag{E94.11}
\]
For the last value one removes the factor three
from \(273\); its remaining unit is \(91\equiv10\),
whose inverse is \(19\) modulo27. The norm identity
is the resultant
\(\operatorname{Res}_Z(Z^6+Z^3+1,1+a_0Z)
 =a_0^6-a_0^3+1\); multiplication of the exponent
of \(\xi\) by any unit modulo9 does not change it.

The actual ninth-radical character has
\(\chi_v(\xi)=5\), a unit modulo9, and
\(\chi_v(\beta_1)=0\). The former follows from
the normalized E.76 local character and the
reciprocity calculation below. The latter follows
from \(N(\beta_3)=\beta_1\). Thus the two characters
are independent on local units. Their joint inertia
image is \(C_9\oplus C_3\), of order27. The
compositum \(E/F\) is totally ramified at these
three primes, so \(E/B_3\) is totally ramified there.
All ramification is in S. This proves (E94.2).

The displayed substitutions in (E94.3) respect
\(\eta^3=\xi\) and \(\beta_3^9=\beta_1\).
Their \(27\cdot9\) distinct actions are
\(G^iH^j\), with multiplication
\[
 (i,j)(k,l)=(i+4^j k,\ j+l)
       \quad\hbox{in }\mathbb Z/27\rtimes\mathbb Z/9.
 \tag{E94.12}
\]
Their number equals \([E:K]\). This proves the
asserted full Galois group and every restriction
in (E94.3). In particular \(G^3,H\) preserve
\(B_3\), while the explicitly different conjugate
of \(\langle H^3\rangle\) proves that \(G\)
does not.

For either ordinary or S-classes, let \(\mathcal H\)
be the corresponding Hilbert three-class field
of \(B_3\). Its intersection with the cyclic cubic
extension \(E/B_3\) is \(B_3\), because that
extension ramifies at \(Q_i\).
For any class in \(\operatorname{Gal}(\mathcal H/B_3)\),
Chebotarev applied to \(E\mathcal H/B_3\) gives a
prime whose Frobenius is that class on \(\mathcal H\)
and identity on \(E\). A degree-one prime above it
has ideal norm representing the given class.
This proves actual norm surjectivity. The precise
source anchors are the norm/Frobenius square,
the Hilbert-class-field correspondence at modulus1,
and Chebotarev
\cite[V,Proposition3.3,Theorem3.6 and Example3.9;
VIII,Theorem7.4]{E_MilneCFT}.

\emph{The exact joint local norm kernels.}
Use the E.76 bases for the three local unit
groups modulo ninth powers. Each is \(C_9^7\).
The three unscaled radical rows and the common
cyclotomic row are
\[
 \begin{pmatrix}
 1&1&5&2&8&3&6\\
 1&1&2&8&5&3&6\\
 1&1&8&5&2&3&6
 \end{pmatrix}\pmod9,\qquad
       (0,0,1,1,1,0,0)\pmod3 .
 \tag{E94.13}
\]
The actual radical rows at these primes are five
times the rows displayed modulo9. For every one
of the three primes the common integer kernel
has the same basis
\[
 K_{\rm loc}=
 \begin{pmatrix}
 9&8&3&6&0&6&3\\
 0&1&0&0&0&0&0\\
 0&0&3&2&2&0&0\\
 0&0&0&1&0&0&0\\
 0&0&0&0&1&0&0\\
 0&0&0&0&0&1&0\\
 0&0&0&0&0&0&1
 \end{pmatrix}.
 \tag{E94.14}
\]
The determinant is27, and the exact Smith quotient
is \(C_9\oplus C_3\).

Here is a literal construction of the new character,
including the exactness test. In
\(F_v[Y]/(Y^3-\xi)\) use the actual norms of \(Y\)
and \(1+(1-Y)^k\). A norm of
\(r+qY+pY^2\) is
\[
 r^3+\xi q^3+\xi^2p^3-3\xi pqr .
 \tag{E94.15}
\]
The first eleven recorded norms, including
\(N(Y)=\xi\), span a six-dimensional subspace
modulo cubes in each of the seven-dimensional
local unit groups. Every column is an actual
norm. Local class field theory gives norm index
three for this totally ramified cubic extension.
Thus these columns give its entire unit norm
kernel. The unique annihilating row is normalized
by its nonzero value on \(\beta_1\), giving
(E94.13) and (E94.11). No unsuccessful bounded
search is used to assert a norm obstruction.

At the primes over seventeen the radical character
on a unit \(u\) is
\(-\log_\xi(\bar u^{32})\) modulo9; the cyclotomic
character is twice the valuation modulo3. On the
fifteen specified S-unit generators of \(F\),
the program evaluates all twelve characters.
Their exact reciprocity rows determine the scale
five in (E94.13). With the actual local generator
choices in the run, the six scaling coefficients
for the radical rows are
\((5,5,5,8,5,4)\). Both sums of actual local
characters vanish on every global S-unit.

The combined character image on all S-units has
order \(3^{13}\); on ordinary units it has order
\(3^7\). On the kernel of the six radical
characters modulo9, the six cyclotomic
characters have rank three. These are integer
kernel/image calculations, with the full
fifteen-column matrices and all Smith maps
retained. The local ideal modulus is
\(Q_1^{16}Q_2^{16}Q_3^{16}T_1T_2T_3\).
It suffices: at a wild prime, logarithm sends
the sixteenth unit filtration inside ninth
powers, since division by nine leaves
filtration greater than \(6/(3-1)=3\).

\emph{The ray-class construction and its maximality.}
Let \(\mathcal R\) be the three-primary ray class
group for that modulus. At each \(Q_i\), impose
the local unit relations (E94.14); at each \(T_i\),
impose ninth powers of units. Chinese remainder
representatives give actual principal-ideal
columns in the ray class group. Their quotient
is \(\Gamma_E=\operatorname{Gal}(\mathcal G_E/F)\).
Indeed an abelian extension over \(F\) becomes
unramified over \(E\) exactly when its local
Artin maps kill \(N_{E_w/F_v}\mathcal O_{E_w}^\times\).
The relations just imposed are those subgroups:
both characters occur at \(Q_i\), and the
unramified cubic cyclotomic factor contributes
no unit condition at \(T_i\).
The local norm/Artin compatibility follows from
the proved restriction/corestriction pairing
\cite[III,Theorem3.1]{E_MilneCFT}; the global
ray-class existence and reciprocity identify
the corresponding field
\cite[V,Theorems3.5--3.6]{E_MilneCFT}.

This is the maximal such three-extension. Away
from S no ramification can disappear under the
unramified extension \(E/F\), so every such
abelian extension is included in the chosen ray
field. At S the imposed subgroups are the full
local unit norm groups, not merely subgroups
generated by some found elements.
The exponent bound used in the computation is
also exact: the remaining local unit quotients
have exponent nine and \(\operatorname{Cl}(F)\{3\}\)
has exponent three. Hence the entire ray quotient
has exponent dividing27. Taking the exponent-27
ray quotient loses no part of \(\Gamma_E\).

On a prime \(\mathfrak l\nmid3\cdot17\) its
two restriction coordinates to \(E/F\) are
computed from
\[
 \bar\beta_1^{(N\mathfrak l-1)/9}
       =\bar\xi^{\,\chi(\mathfrak l)},\qquad
 \bar\xi^{(N\mathfrak l-1)/3}
       =\bar\zeta^{\,\psi(\mathfrak l)} .
 \tag{E94.16}
\]
The certificate factors actual ray basis ideals
into primes and evaluates these two equalities,
including their integer exponents. The kernel
has index27. If \(L_{\rm ord}\) is the full
ray relation lattice, and \(K\) is this
integer character-kernel lattice, the two
presentations are
\[
 \Gamma_E=\mathbb Z^{16}/L_{\rm ord},\qquad
 \mathcal A=K/L_{\rm ord}.
 \tag{E94.17}
\]
The exported bases satisfy \(KP=L_{\rm ord}\);
their Smith factors are respectively
\((9,9,9,9,3,3)\) and \((9,9,9,3)\).
Every Smith left and right matrix has determinant
\(\pm1\), and its inverse is retained.

For complete splitting over \(E\), one also
imposes the full local norm uniformizers. They
can be chosen in \(F_v^\times\) as
\[
 u_{Q_i}=(1+\beta_1)\beta_1^{-2},\qquad
 u_{T_i}=\beta_1^3 .
 \tag{E94.18}
\]
The first has valuation one and both character
values zero by (E94.11) and \(N(1+\beta_3)=1+\beta_1\).
The second has valuation three and both character
values zero. Thus these elements, together with
the local unit norm groups, generate the full
local norm groups. A Chinese remainder element
\(z\) agreeing with \(u_v\) at \(v\), and with
one at the other S-primes, gives the prime-to-S
ideal \(I=(z)/v^{v(u_v)}\). Its negative ray
class is the local Artin class of \(u_v\).
The program constructs all six such ideals
and checks their S-valuations exactly.

Adjoining these six columns to \(L_{\rm ord}\)
gives a lattice \(L_S\subset K\). Its quotient
over \(F\) has Smith factors \((9,9,3,3)\);
\(K/L_S\) has factors \((9,3)\).
This proves the second row of (E94.4).
The inclusion \(L_{\rm ord}\subset L_S\) gives
the actual quotient matrix (E94.6), and its
six columns, written in the first relative
Smith basis, are exactly (E94.7).
They are also the Frobenii of \(Q_i^E,T_i^E\):
local norm compatibility takes a uniformizer
upstairs to a local norm uniformizer downstairs.
Their generated subgroup has order81, equal
to the kernel order of (E94.6).

\emph{The original-genus norm and the genuine extension map.}
The local unit norms from \(E/F\) lie in those
from \(B_3/F\). Hence the original ninth-genus
field is a subfield of \(\mathcal G_E\).
The program constructs its restriction matrix
by evaluating each actual new ray basis ideal
in the original ray class group. Equality of
cyclic-factor lists is not used to identify
the two ray bases.
Intersecting with the two Kummer kernels gives
the onto matrix \(N\) in (E94.8), in exactly
the simple coordinates of E.91. For an upper
prime of residue degree \(f\), restriction of
its Frobenius is the \(f\)th power of the
lower Frobenius. Thus this restriction is
the actual ideal norm map.

Let \(j\) be ordinary ideal extension from
\(B_3\) to \(E\), followed by (E94.5).
The lift \(G^3\) of \(g_2\) lies in
\(\operatorname{Gal}(E/F)\), so it acts
identically on the abelian-over-\(F\) genus
group. Equivariance therefore proves
\(j((g_2-1)c)=0\) for every actual class \(c\).
E.91 identifies that subgroup with the
complete kernel \(W\) of its original genus
map. This proves (E94.9), and gives a
well-defined map \(J:A_9\to\mathcal A\).

The norm from \(\operatorname{Cl}(E)\{3\}\)
onto \(\operatorname{Cl}(B_3)\{3\}\) is onto,
as proved above. On the new genus quotient
the deck group \(\langle H^3\rangle\) acts
identically. The actual ideal identities give
\(JN=1+H^3+H^6=3\) and \(NJ=3\).
The full Smith calculation of \(\ker N\)
gives \(C_3^2\); hence three times any lift
under \(N\) is independent of the chosen
lift. Applying this to the three specified
old basis classes gives exactly the displayed
matrix \(J\). This proves that the finite
matrix is the actual extension comparison.

The six prime columns provide independent
checks of the coefficients:
\[
 NC_E=C_{\rm old}\operatorname{diag}(1,1,1,3,3,3),
 \qquad
 JC_{\rm old}=C_E\operatorname{diag}(3,3,3,1,1,1).
 \tag{E94.19}
\]
Here \(C_{\rm old}\) is the E.91 matrix in
its \((y,z,t)\) coordinates. Thus the new
prime norms do not silently identify the
inert and ramified prolongations.

\emph{Normalizer action and the exact noncyclic comparison.}
Both \(G,H\) preserve \(F\) and normalize
the defining norm subgroups. Their actions
on \(F\) are the literal substitutions
\(a\mapsto\zeta a,\xi\mapsto\xi\) and
\(a\mapsto\zeta^{-1}a,\xi\mapsto\xi^4\).
The actual ray Galois matrices preserve
\(L_{\rm ord},L_S,K\), and induce on
\(\mathcal A\) the matrices
\[
 G_{\mathcal A}=
 \begin{pmatrix}7&1&0&0\\0&7&6&0\\6&4&7&0\\0&2&1&1\end{pmatrix},
 \qquad
 H_{\mathcal A}=
 \begin{pmatrix}1&1&3&6\\0&7&3&0\\6&5&7&0\\1&1&1&1\end{pmatrix}.
 \tag{E94.20}
\]
They commute and their cubes are identity,
with the indicated row moduli. Taking the
quotient (E94.6) gives precisely (E94.10).

The S-class norm to \(F\), evaluated on the
same actual ray basis ideals, is its certified
second class coordinate modulo three. E.86
identifies the norm
\(\operatorname{Cl}_S(B_3)\to\operatorname{Cl}_S(F)\)
with the identity in the coherent prime-\(73\)
bases. Therefore this character is exactly
the norm to the original \(C_3\), not just
some nonzero character of that order.
Its value in the two relative Smith coordinates
is \((-1,0)\), proving (E94.10).

The source of (E94.5) modulo \(H^3-1\)
still surjects onto \(\mathcal A_S\), since
\(H^3\) belongs to the abelian genus base
group. A cyclic
\(\mathbb Z_3[\langle H^3\rangle]\)-module
has cyclic \(\mathbb Z_3\)-coinvariants:
the image of a single module generator
generates the quotient. Its quotients
are cyclic. The explicitly proved quotient
\(C_9\oplus C_3\) has two independent
modulo-three generators. This proves the
noncyclic assertion through the actual
surjective map, and identifies the exact
premise that fails in transporting the
earlier cyclic-module computation to \(E\).

Finally the Python verifier checks ten full
Smith systems, both ideal norm/extension
identities, all \(2187\) relative ordinary
coordinates, all \(27\) relative S coordinates,
and the \(243\) normalizer group elements.
The six-stage portable rebuild begins with
the certified \(F\)-BNF, reconstructs every
new local and ray object, and retains its
actual stdout and binary mathematical objects.
The remaining values excluded from the
statement are not inferred from the genus
quotients or from the two retained failed
local spanning attempts.
\end{proof}
